\lab{11}{Сжатие индекса}
\task
Применить сжатие к координатным блокам, описать битовое представление,
исследовать размеры и скорость прохода и поиска. Проверить сохранение
результатов и качества при преобразовании без потерь.

\method
Выбран unsigned LEB128: младшие 7 бит байта содержат очередную группу
разрядов числа, старший бит равен 1 при наличии следующего байта.
Группы записываются от младших разрядов к старшим. Для 32-битового числа
достаточно от одного до пяти байтов.
\begin{table}[H]\centering
\begin{tabular}{rll}\toprule
Число & Байты, hex & Биты\\\midrule
0 & \code{00} & \code{00000000}\\
127 & \code{7F} & \code{01111111}\\
128 & \code{80 01} & \code{10000000 00000001}\\
16383 & \code{FF 7F} & \code{11111111 01111111}\\\bottomrule
\end{tabular}
\caption{Кодирование целых LEB128.}
\end{table}
Вместо абсолютных \(docID\) записываются разности с предыдущим
документом. Позиции также кодируются разностями, но предыдущая позиция
сбрасывается в начале каждого документа. Частота вхождений остаётся
обычным беззнаковым числом, закодированным LEB128. Нулевой первый
идентификатор и нулевая первая позиция допустимы.

В raw-варианте хранятся абсолютные \(docID\) и позиции по 4 байта.
Таким образом, сравниваются практические представления raw и
«разности + VB», а не только замена кодека при одинаковых числах.
Для \(P\) постингов и \(T\) позиций raw-координаты занимают
\(4(2P+T)\) байт, что позволяет независимо сверить размер секции.

\begin{lstlisting}[language=C++]
do {
    unsigned char byte = value & 127;
    value >>= 7;
    if (value) byte |= 128;
    require(fwrite(&byte, 1, 1, file) == 1, "Write error");
} while (value);
\end{lstlisting}
Декодер контролирует конец секции, переполнение пятого байта
и неканоническое представление с лишним нулевым старшим блоком.

\journal
Чтение и запись целых вынесены в общий кодек, формат помечен флагом
сжатия. При построении сохраняются те же термы, частоты и метаданные.
В эксперименте прыжки отключены, чтобы изменение их таблицы не
влияло на сравнение размеров. Поисковые запросы выполняются после
одного прогревочного прохода, затем повторяются семь раз.
Построение и измерение разделены: измеритель читает готовые raw-
и VB-индексы и не требует исходного корпуса. Повторная подготовка
raw-индекса при существующем файле отклоняется без удаления
индекса и других результатов эксперимента.

\testing
Кодирование и декодирование проверяются на 0, границах 7-битовых
групп, \code{UINT32_MAX} и 1000 воспроизводимых псевдослучайных числах.
Проверяются некорректные байты и обрыв данных. Поиск по raw, VB
и индексу с многоблочной сборкой должен возвращать один эталонный
список. Четыре теста повторно пройдены. В полном эксперименте
дополнительно сравниваются статистика и контрольная сумма результатов.

\results
\begin{table}[H]\centering
\begin{tabular}{lrr}\toprule
Показатель & raw & VB\\\midrule
Полный индекс, байт & 16\,694\,523\,464 & 6\,510\,260\,552\\
Координатная часть, байт & 15\,056\,508\,824 & 4\,872\,245\,912\\
Прыжки, байт & 0 & 0\\
Среднее время запроса, мс & 42,546 & 48,289\\
Медианное время, мс & 17,787 & 20,540\\
Прочитано постингов & 39\,689\,783 & 39\,689\,783\\\bottomrule
\end{tabular}
\caption{Одинаковые 30 запросов, семь повторов после прогрева. Число прочитанных постингов суммарное по измеряемым повторам.}
\end{table}
Полный файл сократился до 38,996\% исходного размера (экономия
61,004\%, сжатие в 2,56 раза). Координатная часть сократилась примерно
на 67,64\%, в 3,09 раза. Среднее время выросло на 13,50\%.
Контрольные суммы результатов поиска в обоих случаях совпадают.
Числа документов, термов, постингов и позиций совпадают.

\begin{figure}[H]\centering
\begin{tikzpicture}
\begin{axis}[ybar,width=0.8\textwidth,height=6cm,
ylabel={Размер, GiB},symbolic x coords={raw,VB},xtick=data,
ymin=0,bar width=24pt,legend style={at={(0.5,1.02)},anchor=south,legend columns=2}]
\addplot coordinates {(raw,15.548),(VB,6.063)};
\addplot coordinates {(raw,14.022),(VB,4.538)};
\legend{Полный индекс,Координаты}
\end{axis}
\end{tikzpicture}
\caption{Изменение размера. Координаты входят в полный индекс; столбцы не суммируются.}
\end{figure}
\begin{table}[H]\centering\small
\begin{tabular}{lrrrrr}\toprule
Группа & Термов & raw, MiB & VB, MiB & raw, с & VB, с\\\midrule
Все & 20\,430\,767 & 14359,01 & 4646,54 & 165,469 & 176,139\\
Редкие & 17\,693\,418 & 340,32 & 145,71 & 18,408 & 18,034\\
Средние & 2\,302\,690 & 248,61 & 102,24 & 9,033 & 9,217\\
Частые & 434\,659 & 13770,08 & 4398,59 & 148,605 & 156,828\\
\bottomrule\end{tabular}
\caption{Объём координат и время последовательного прохода по группам документной частоты.}
\end{table}

Редкими считаются термы с документной частотой 1--2,
средними с частотой 3--19, частыми с частотой от 20.
Для полного прохода время выросло с 165,469 до 176,139 секунды,
на 6,45\%. У редких термов проход немного ускорился, на 2,03\%,
а у средних и частых замедлился на 2,04\% и 5,53\% соответственно.
Частые термы сжимаются плотнее: координатная часть составляет
31,94\% raw-размера против 42,82\% для редких термов.
Контрольные суммы, число постингов и число позиций совпадают
во всех группах. Время VB-сборки отдельно не измерено;
поисковое время не заменяет время построения.

\conclusions
Сжатие уменьшило объём существенно, но на данной нагрузке стоимость
декодирования оказалась выше выигрыша от уменьшения ввода-вывода.
Ускорение нельзя выводить только из размера файла. Совпадение
контрольных сумм и тестовых ответов подтверждает сохранение выдачи
на проверенных запросах. При одинаковых оценках релевантности
неизменная выдача сохраняет и значения поисковых метрик.
Возможны блочные кодеки и буферизованное
декодирование, но их выигрыш следует измерять отдельно.
